\documentclass[12pt,a4paper]{article}
\usepackage[utf8]{inputenc}
\usepackage[T1]{fontenc}
\usepackage{geometry}
\usepackage{setspace}
\usepackage{array}
\usepackage{booktabs}
\usepackage{multirow}
\usepackage{longtable}
\usepackage{enumitem}
\usepackage{float}

\geometry{left=2.5cm, right=2.5cm, top=2.5cm, bottom=2.5cm}
\setstretch{1.15}

% Comando para gerar um formulário de extração de dados numerado (em branco)
\newcommand{\formularioExtracaoDados}[1]{%
  \subsection*{#1º)}%
  \noindent \textbf{Título do Artigo:} \\
  \textbf{Autores:} \\
  \textbf{Data da Publicação:} \\
  \textbf{Veículo de Publicação:} \\
  \textbf{Fonte:} \\
  \textbf{Abstract:} \\
  \textbf{Resumo:} \\
  \textbf{Estudo:} \\
  \textbf{Data de execução:} \\
  \textbf{Local:} \\
  \textbf{Tipo:} \\
  \textbf{Descrição:} \\
  \textbf{Hipóteses avaliadas:} \\
  \textbf{Variáveis independentes:} \\
  \textbf{Variáveis dependentes:} \\
  \textbf{Participantes:} \\
  \textbf{Material:} \\
  \textbf{Planejamento do estudo:} \\
  \textbf{Ameaças à validade:} \\
  \textbf{Resultados:} \\
  \textbf{Comentários adicionais:} \\
  \textbf{Referências relevantes:} \\
  \textbf{Justificativa:}%
}

% Macro auxiliar: só imprime o campo se o valor não estiver vazio
\newcommand{\campoForm}[2]{%
  \ifx\empty#2\empty
  \else
    \noindent \textbf{#1} #2\\%
  \fi
}

% Comando para gerar um formulário de extração de dados já preenchido
% ATENÇÃO: o LaTeX só permite até 9 parâmetros por comando.
% Uso: \formularioExtracaoDadosPreenchido{n}{Título}{Autores}{Data}{Veículo}{Fonte}{Abstract}{Resumo}{Estudo}
\newcommand{\formularioExtracaoDadosPreenchido}[9]{%
  \subsection*{#1º)}%
  \campoForm{Título do Artigo:}{#2}%
  \campoForm{Autores:}{#3}%
  \campoForm{Data da Publicação:}{#4}%
  \campoForm{Veículo de Publicação:}{#5}%
  \campoForm{Fonte:}{#6}%
  \campoForm{Abstract:}{#7}%
  \campoForm{Resumo:}{#8}%
  \campoForm{Estudo:}{#9}%
}

\begin{document}

\begin{center}
    {\Large \textbf{PROTOCOLO DE REVISÃO}}
\end{center}
\vspace{1cm}

\noindent \textbf{Pesquisador(a):} Noedson Romario Vieida da Silva

\vspace{0.5cm}
\noindent \textbf{Tema:} Agente Inteligente Para Minotoramento da Quaiade do Ar em Smart Campus

\vspace{0.5cm}
\noindent \textbf{Objetivos:} Identificar estudos que propóem arquiteturas
IoT (LoRaWan, FIWARE, Chirpstack, Node-Red) para monitoramento da qualidade
do ar em ambientes de Smart Campus, com foco em agentes inteligentes para 
análise e previsão da qualidade do ar;


Identificar os principais desafios e soluções propostas para a implementação 
dessa arquitetura, incluindo aspectos de coleta de dados, processamento em 
tempo real e integração com sistemas de alerta;

Identificar sistemas que integram sistemas operacionais de tempo real (RTOS) 
em nós sensores IoT para garantir respostas determinísticas;

\vspace{0.5cm}
\noindent \textbf{Formulação da(s) pergunta(s) da revisão:}
\begin{itemize}[label={\textbullet}, nosep]
    \item Quais arquiteturas IoT (sensores, gateways LoRaWAN, middleware, nuvem) 
          têm sido usadas para monitorar qualidade do ar em smart campuses?
    \item Quais algoritmos de IA/ML são aplicados para predição de curto prazo 
          (próximas horas) de poluentes em ambientes internos/externos?
    \item Existem soluções que integram RTOS em nós sensores para garantir respostas 
          determinísticas em monitoramento de qualidade do ar?
    \item Quais mecanismos de resposta automatizadas são descritos na literatura?
\end{itemize}

\vspace{0.5cm}
\noindent \textbf{Fonte(s):} IEEE Xplore, e Google Scholar, Periódico CAPES, Scopus, Web of Science

\vspace{0.5cm}
\noindent \textbf{Data/período da Busca:} 2021 - 2026

\vspace{0.5cm}
\noindent \textbf{Intervalo de tempo da busca:} 5 anos

\newpage
\vspace{0.5cm}
\noindent \textbf{Palavras-chave:}
\vspace{0.3cm}

\begin{table}[H]
\centering
\begin{tabular}{|p{4 cm}|p{4 cm}|p{4 cm}|p{4 cm}|}
\hline
\textbf{Palavra-chave em Português} & \textbf{Sinônimos em Português} & \textbf{Palavra-chave em Inglês}& \textbf{Sinonimos em Inglês}\\
\hline
Qualidade do ar & Condição do ar ambiente, Nível de poluição do ar & Air quality & Ambient air quality, Air Pollution, Atmospheric quality\\
\hline
Smart campus & Campus inteligente, campus conectado & Smart campus & Intelligent campus, Connected campus, Digital campus\\
\hline
IoT & Internet das Coisas, Dispositivos conectados, redes de sensores & IoT & Internet of Things, IoT Systems, Smart Devices\\
\hline
LoRaWAN & Rede LPWAN de longo alcance, Infraestrutura LoRaWAN para IoT & LoRaWAN & LoRaWAN Network, LoRaWAN Infrastructure\\
\hline
FIWARE & Plataforma de gerenciamento de contexto FIWARE, middleware FIWARE para IoT & FIWARE & FIWARE Platform, IoT Middleware\\
\hline
Predição de poluentes & Predição de curto prazo de poluentes & Short-term pollutant prediction & Short-term Air Quality Prediction\\
\hline
\end{tabular}
\end{table}

\vspace{0.5cm}
\noindent \textbf{String de busca (inglês):}
\begin{verbatim}
- (("air quality" OR "air pollution" OR "indoor air quality") AND 
   ("smart campus" OR "university campus" OR "higher education") AND 
   ("IoT" OR "LoRaWAN" OR "FIWARE" OR "ChirpStack") AND 
   ("predictive" OR "forecast" OR "anomaly detection") AND 
   ("real-time" OR "RTOS" OR "FreeRTOS" OR "deterministic")) 
  NOT ("wearable" OR "smart home" NOT campus)
\end{verbatim}

\vspace{0.3cm}
\noindent \textbf{String de busca (português):}
\begin{verbatim}
- (("qualidade do ar" OR "poluição do ar" OR "qualidade do ar interior") AND 
   ("campus inteligente" OR "campus universitário" OR "ensino superior") AND 
   ("IoT" OR "LoRaWAN" OR "FIWARE" OR "ChirpStack") AND 
   ("preditivo" OR "previsão" OR "detecção de anomalias") AND 
   ("tempo real" OR "RTOS" OR "FreeRTOS" OR "determinístico")) 
  NOT ("vestível" OR "casa inteligente" NOT campus)
\end{verbatim}



\newpage
\section*{Critérios de Inclusão e Exclusão}

\begin{table}[ht]
\centering
\begin{tabular}{|p{10cm}|p{5cm}|}
\hline
\textbf{Critérios} & \textbf{Tipo} \\
\hline
 & \\
\hline
 & \\
\hline
 & \\
\hline
 & \\
\hline
 & \\
\hline
 & \\
\hline
\end{tabular}
\end{table}

\noindent \textbf{Justificativa:}

\vspace{0.5cm}
\section*{Critérios de Qualidade}

\begin{table}[ht]
\centering
\begin{tabular}{|p{10cm}|p{4cm}|}
\hline
\textbf{Critérios} & \textbf{Nota (De 0 a 2)} \\
\hline
1. & \\
\hline
2. & \\
\hline
3. & \\
\hline
4. & \\
\hline
5. & \\
\hline
\end{tabular}
\end{table}

\noindent \textbf{CLASSIFICAÇÃO = Somatória das notas dos critérios (N)} \\
\textbf{OBS:} \\
N $\ge$ 85\% (Excelente) \\
65\% $\le$ N $\le$ 85\% (Muito Boa) \\
45\% $\le$ N $\le$ 65\% (Boa) \\
25\% $\le$ N $\le$ 45\% (Média) \\
N $<$ 25\% (Baixa)

\section*{Lista dos artigos encontrados}
\subsection*{Lista dos artigos incluídos}

\begin{table}[ht]
\centering
\begin{tabular}{|c|p{6cm}|c|c|c|}
\hline
\textbf{Nº} & \textbf{Título do artigo} & \textbf{Autores} & \textbf{Publicação} & \textbf{Veículo} \\
\hline
1 & \# & \#\# & & \\
\hline
2 & & \#\# & & \\
\hline
3 & \# & & & \\
\hline
4 & & & & \\
\hline
5 & & \#\# & & \\
\hline
6 & & & & \\
\hline
7 & & & & \\
\hline
8 & & & & \\
\hline
9 & & & & \\
\hline
10 & & & & \\
\hline
11 & & & & \\
\hline
12 & & & & \\
\hline
13 & & & & \\
\hline
14 & & & & \\
\hline
15 & & & & \\
\hline
16 & & & & \\
\hline
\end{tabular}
\end{table}

\subsection*{Lista de artigos excluídos}

\begin{table}[h]
\centering
\begin{tabular}{|c|p{6cm}|c|c|c|}
\hline
\textbf{Nº} & \textbf{Título do artigo} & \textbf{Autores} & \textbf{Publicação} & \textbf{Veículo} \\
\hline
1 & & & & \\
\hline
2 & & \#\# & & \\
\hline
3 & & & & \\
\hline
4 & & & & \\
\hline
5 & \# & & & \\
\hline
\end{tabular}
\end{table}

\newpage
\section*{FORMULÁRIO DE EXTRAÇÃO DE DADOS}

% Exemplo: formulário 1 já preenchido
% Substitua os textos entre chaves pelos dados do seu artigo
% \formularioExtracaoDadosPreenchido{1}{Título}{Autores}{Data}{Veículo}{Fonte}{Abstract}{Resumo}{Estudo}

% Exemplo real preenchido (campos extras podem ser deixados vazios com {})
% \formularioExtracaoDadosPreenchido
%   {1}% número do formulário
%   {Título do artigo}% título
%   {Noedson Romario de Oliveira, Gustavo Pinto, Marco Tulio Valente}% autores
%   {2020}% data da publicação
%   {IEEE Transactions on Software Engineering}% veículo
%   {IEEE Xplore}% fonte
%   {Abstract do artigo...}% abstract
%   {Resumo do artigo...}% resumo
%   {Tipo do estudo...}% estudo

% Formulário em branco (para preencher à mão ou no PDF)
\newpage
\section*{ANÁLISE DOS RESULTADOS}

\noindent Agrupamento, comparação e discussão crítica dos trabalhos relacionados

\vspace{0.3cm}
\noindent - Tabulação de resultados (Dados quantitativos)

\vspace{0.3cm}
\noindent \hspace{1cm} - \textbf{O que já foi publicado?}

\vspace{0.3cm}
\noindent \hspace{1cm} - \textbf{Quais as teorias utilizadas por estudo?}

\vspace{0.3cm}
\noindent \hspace{1cm} - \textbf{O que é conhecido até recentemente sobre o tema?}

\vspace{0.3cm}
\noindent - \textbf{Indicadores/Métricas/Critérios e seus valores}

\vspace{0.3cm}
\noindent - \textbf{Comparações críticas}

\vspace{0.3cm}
\noindent - \textbf{Lacunas existentes}

\newpage
\section*{REFERÊNCIAS}

\vspace{5cm}
\noindent

\end{document}